\section{Transcendence and the Aristocratic Principle}

\begin{quotex}
The cause of all wars and revolutions---in a word, of all violence---is always the same: negation of hierarchy.
\flright{\textsc{Valentin Tomberg}}
\end{quotex}


In the essay \emph{Transcendence and the Aristocratic Principle}, published in Aristokratia III\footnote{\url{http://manticore.press/aristokratia/aristokratia-iii/}}, \textbf{Edwin Dyga} gives us an excellent overview of traditional reactionary thinkers, that is, those of the ``Old Right''. Specifically, he identifies the ``Throne and Altar'' as the essential criteria for civilization or, as we would say, the political power and spiritual authority in their proper relationship. He also touches on some so-called New Right thinkers, but, as far as I understand his point, they are defective in dealing with these criteria in a new way.

The opposite, then, of the Throne and Altar paradigm is ``Revolution''. Mr. Dyga then identifies two propositions resulting from the rejection of the paradigm:

\begin{itemize}


\item The spirit of the Revolution involves a rejection of an order symbolized by ``Throne and Altar''

\item The ideology of the Revolution are demotic and materialist in essence
\end{itemize}


The first proposition is axiomatic, but the second requires an explanation. The two terms, demotic and materialist, are not connected by happenstance, but rather by necessity. \textbf{Valentin Tomberg} explains why:

\begin{quotex}
Without an Emperor, there will be, sooner or later, no more kings. When there are no kings, there will be, sooner or later, no more nobility. When there is no more nobility, there will be, sooner or later, no more bourgeoisie or peasants. This is how one arrives at the dictatorship of the proletariat, the class hostile to the hierarchical principle which, however, is the reflection of divine order. This is why the proletariat professes atheism.
\end{quotex}


The demotic essence, therefore, is atheistic, anti-hierarchical, materialistic. Since there are New Right movements that themselves are atheistic, anti-hierarchical, and materialistic, this just shows the power the modernity holds over some minds. Such movements as typically ``identitarian'', often on a biologically racial basis, which then becomes the defining paradigm rather than Throne and Altar. That is because such movements are themselves demotic and therefore unwilling to acknowledge or unable to recognize the hierarchy within the group itself.

Since the reader is urged to read the essay himself, we will be content to include a few traditional principles to aid in further developments along these lines.


\paragraph{Real and Ideal Relations}

The philosopher \textbf{Timothy Sprigge} has given us the useful distinction between real and ideal relations. Real relations are related materially or historically; that is, there is a direction relationship on the horizontal plane. An ideal relationship is vertical; two ideas are not historically connected, yet they derive their similarity from the same transcendent principles.

So, for example, the Old Right derived from opposition to the French revolutionary spirit. Old Right thinkers are then Joseph de Maistre, de Bonald, Donoso Cortes as well as the more contemporary thinkers Mr. Dyga mentions. They are in a real relationship.

The New Right, on the other hand, is not in that tradition. Rather, it uses modern and post-modern thought to define a new understanding of the right that is not tied to the restoration of Throne and Altar. To the extent it that it defines a homologous alternative, it is in an ideal relationship to the Old Right. Specifically, it needs to define an alternative political power arrangement and a legitimate spiritual authority. Failing that, it devolves into a pale imitation of modernism.

An example is the Old and New Lefts. The Old Left saw the rejection of order solely in economic terms: the proletariat would take over the means of production. The New Left likewise sought a rejection of hierarchical order. However, they defined this first in racial terms, followed later by sexual relationships and so on. So, historically they arose from different motivations. However, since they are both revolutionary, demotic, and materialistic movements, they are in an ideal relationship to each other.

\paragraph{Homology}


\textbf{Julius Evola} uses the term ``homology'' in a similar sense. For example, he recognizes three main civilizations descending from the Hyperboreans: the Vedic, the Greco-Roman pagan, and the Nordic-Roman medieval civilizations. In them, the same patterns repeat, i.e., they are homologous. For example, they are socially hierarchical with a spiritual, warrior, and producing classes. Georges Dumezil recognized this as common to Indo-European civilizations.

There is no need to belabor this here, since we have provided sufficient examples already. The point is that any new right movement needs to explain how the civilization it envisions is homologous to other traditional civilizations.

\paragraph{The Illusion of the Modern World}


Mr. Dyga points out that

\begin{quotex}
In this critique of the modernist world, Rene Guenon accordingly states that the foundational assumptions of the modern world are a contradiction of the cosmic order.
\end{quotex}


The real impact of this claim is seldom noted. First of all, the corollary is that the traditional world is based on the cosmic order. Hence, it is incumbent on all rebels and soi-disant ``aristocrats of the soul'' to understand what that order is and how it manifests in a particular place at a particular time.

The second corollary is that the modern world is unreal, as it is based on an illusion. Just as a chemist cannot deny the law that water boils at 100?, neither can a political scientist deny the cosmic order. Hence, the dispute between the modernist and the traditionalist is not an intellectual battle of ideas, but rather the difference between illusion and truth.

\paragraph{Degeneration or the Third Dimension}


The real battle then is spiritual, not intellectual. Specifically, it is ultimately futile to try to trace the degeneration of civilization as a logical sequence of ideas, with one proposition leading to the next. This is just a form of historicism.

In fact, the degeneration begins in the people, or demos, themselves, not in the ideas they hold. Specifically, certain ideas can only take hold in minds that are already degenerate. The recognition of the spiritual or demonic origin of certain streams of thought constitutes the third dimension of history. Mr. Dyga also recognizes this, if not always consistently:

\begin{quotex}
It is not the details that are of primarily interest to us but the animating force behind them, the current which emanates from an ephemeral realm yet whose force shapes the nature of temporal phenomena.
\end{quotex}


\paragraph{Quo Vadis}


After all that theorizing, the temptation is to create a political program of some sort. How, then, can you arrive at the opposite of a revolution? It cannot be demotic (i.e., a mass movement); hence it requires an elite who understand the traditional principles and axioms. It cannot be atheistic or materialistic.

In his dialog with the New Right, Mr. Dyga walks a tightrope. If identity is primary, as the New Right presumes, then any spiritual movement is merely instrumental: it either supports or opposes the identity in question. Mt. Dyga tries, contrarily, to demonstrate that the hierarchical spiritual authority will create identity as a consequence. The principle of subsidiarity guarantees this. Historically, of course, this has been true. Christian Europe simultaneously repelled outsiders on the one hand, while simultaneously respected a Europe of 100 flags.

It was only the widespread rejection of Christianity, particularly in its most Traditional manifestation (specifically the Nordic-Roman Medieval Church), that threatened that understanding. To his credit, Mr. Dyga refutes the worst impulse of the New Right in a reasoned and calm manner. He often appeals to Julius Evola in rejecting the most absurd claims of zoological racism. Nevertheless, he recognizes that there are differences between peoples, but that they are the result of group spiritual and psychical/cultural factors.

When Peter was fleeing Rome to avoid persecution, he encountered the risen Christ. ``Quo Vadis?'' Peter asked. (``Where are you going?'') Jesus replied, ``I am going to Rome.'' Hearing that, Peter turned around and returned to Rome. That seems to be the response of Mr. Dyga.

\paragraph{Postscript on Tragedy}


Since the essay refers to \textbf{Jean Raspail}'s \emph{Camp of the Saints}, given the timing of current events, it merits a comment. Raspail's book should not be understand as a racist screed, but rather as a satire of the political and religious leaders of France. Curiously, life imitates art, and the events foreseen in the novel are actually coming to pass. There is even the South American pope that Raspail predicted. There are Asian refugees entering Europe by the hundreds of thousands. There are the same horrible scenes, both in the book and in life, of children drowning in the sea.

Yet, it goes on, whatever the cost. And the political and religious leaders speak the same words as though Raspail were feeding them their lines. The Tragedy of Life arises from certain features of the players, called ``hamartia'', whether good, bad, or comic. While the audience can anticipate how the whole thing will play out, the actors are oblivious to the consequences of their decisions. They cannot do otherwise. That is the Tragedy.


\flrightit{Posted on 2015-09-03 by Cologero}

\begin{center}* * *\end{center}

\begin{footnotesize}
\begin{sffamily}
\texttt{Alistair Fraser on 2015-09-04 at 07:14 said:}

From above... ... ... 

``While the audience can anticipate how the whole thing will play out, the actors are oblivious to the consequences of their decisions. They cannot do otherwise. That is the Tragedy.''

And speaking in the spirit of the piece,

The Actors (politicians) and also much of the audience, although that balance is swinging now regarding the audience, are all ``Drunk on idea'' it is a true life drama of the tragedy of accountancy dominated ``opinionated'' human thought patterns played out on the world of mediated news events. These actors and their actions and decisions are a perfect dramatical display of the result of the calcification of human thought , whereby the bottom line is always new short term knee jerk ``ideas'' administered by political experts .

What is required is a spiritual guidance that is superior to these nodding head political monetary worshippers, Unfortunately in these times the official sources of spirituality are also worshippers of the monetary overlord. Somewhere out there , there are minds that could administer spiritual nourishment and guidance , but these people tend to be ridiculed if they ever put one foot on a public pedestal, So who will attempt to administer the ``calcination process'' on the current monetary overlord thought system

\hfill

\texttt{william zeitler on 2015-09-04 at 09:32 said:}

``In the essay Transcendence and the Aristocratic Principle, published in Aristokratia III,... '' contains a link to \url{http://www.aristokratia.info/buy-aristokratia.html} presumably so one can purchase/read the article -- this link is broken. Ed. Apparently there was a merger and a new web site:

\url{http://manticore.press/aristokratia/aristokratia-iii/}

\hfill

\texttt{Thomas Blanchard on 2015-09-04 at 11:18 said:}

In light of recent events in Europe, the foresight of Plinio Corrêa de Oliveira in his 1993 book, Nobility and Analogous Traditional Elites in the Allocutions of Pius XII, is rather impressive:

\url{http://www.tfp.org/tfp-home/books/nobility.html}


``It is not improbable that an armed conflict within the former U.S.S.R. would lead to the involvement of major Western nations, with consequences of apocalyptic dimension. One of these consequences could easily be the migration of entire populations pressed by fear of war and actual famine to Central and Western Europe. This migration could assume a critical character of unpredictable scope.

What effects would this exodus have on nations until recently under Soviet domination, such as those on the Baltic Sea? What effects would it have on other countries, such as Poland, the Czech Republic, Slovakia, Hungary, Rumania, and Bulgaria, about which it would be very daring to affirm that they have entirely escaped the communist yoke?

To complete this panorama, we should consider the possible reaction of the Maghreb in face of a Western Europe enmeshed in problems of this magnitude, as well as developments throughout northern Africa and the profound impact of the immense fundamentalist wave sweeping the peoples of Islam, of which the Maghreb is an integral part. Who can predict with certainty the extremes to which these factors of instability will bring the world, and especially the Christian world?

For the time being, the latter is not engulfed in the triple drama of a seemingly peaceful invasion from the East, a probably less peaceful invasion from Africa, and an eventual worldwide conflagration. However, the fatal outcome of the long revolutionary process whose outline was summarized in the last chapter of this work is already within sight.

This process has advanced relentlessly, from the waning and fall of the Middle Ages to the initial joyful triumphs of the Renaissance; to the religious revolution of Protestantism, which remotely began to foment and prepare the French Revolution and, even more remotely, the Russian Revolution of 1917. So invariably victorious has been its path despite uncountable obstacles that one might consider the power that moved this process invincible and its results definitive.

These results seem definitive indeed if one overlooks the nature of this process. At first glance it seems eminently constructive, since it successively raised three edifices: the Protestant Pseudo-Reformation, the liberal-democratic republic, and the Soviet socialist republic.

The true nature of this process, however, is essentially destructive. It is Destruction itself. It toppled the faltering Middle Ages, the vanishing Ancien Régime, and the apoplectic, frenetic, and turbulent bourgeois world. Under its pressure the former U.S.S.R. lies in ruins---sinister, mysterious, and rotten like a fruit long-since fallen from the branch.

Hic et nunc, is it not true that the milestones of this process are but ruins? And what is the most recent ruin generating but a general confusion that constantly threatens imminent and contradictory catastrophes, which disintegrate before falling upon the world, thus begetting prospects of new catastrophes even more imminent and contradictory. These may vanish in turn, only to give way to new monsters. Or they may become frightful realities, like the migration of Slavic hordes from the East to the West, or Moslem hordes from the South to the North.''

\hfill

\texttt{Mark Citadel on 2015-09-05 at 17:08 said:}

Another excellent article from Gornahoor! Perfectly elucidates the failures of the New Right vis-a-vis the Old Right. No Reactionary should miss this piece, in fact I am compelled to write something regarding it in the coming days.

You pointed out that some New Right groups are atheistic. Does this entire criticism hold true against the neo-Pagan elements of the New Right, or only the atheistic ones?

\hfill

\texttt{Cologero on 2015-09-06 at 00:07 said:}

Mark,

That is a distinction without a difference.

Those ``neo-pagan'' elements are also atheists. They don't recognize the sort of paganism that we've described on this site, e.g., from the sanatana dharma to Plato/Aristotle/Plotinus. Rather, they prefer the vulgar polytheism of the masses, while, of course, not actually believing in the actual existence of those gods and goddesses. Moreover, as we've also pointed out, those pagan gods represent real forces, e.g., Demeter, Aphrodite, etc., not necessarily benign. The pagans were concerned about propitiating one god(ess) or another. That is ignored by the new right neo-pagans (NRNP).

The real motivation is anti-Christian and nothing else, even to the point of denying historical reality and justifying absurdities such as the conflation of liberalism with Christianity.

Things came to a head at this site a couple of years ago when we criticized \textit{Dominique Venner's suicide at Notre Dame}\footnote{\url{http://www.gornahoor.net/?p=6475}}. We lost about 30\% of the page views right after that.

For an example, one of those NRNPs wrote a facebook post praising Hungary's stand against immigration. Of course, he curiously omitted the reason: the prime minister wanted to ``preserve Europe's Christian identity''! The NRNP has no desire for such a preservation, as a matter of fact, just the opposite. Moreover, the least Christian nations of Europe --- Germany, the UK, France, Sweden --- are the nations most in favor of it.

Dominique Venner despised Christian France and wrote that he had wished France had been more Hindu. Contrast him with his contemporary, Jean Raspail, who superficially sound alike, yet with a much different view of things, especially the desirability of Hinduism in France. Go figure.

\hfill

\texttt{Cologero on 2015-09-06 at 00:10 said:}

Another point: I wonder if those European neo-hindu, neo-pagans have sorted out among themselves which caste they would belong to.

\hfill

\texttt{E. DYGA on 2015-09-08 at 09:27 said:}

My gratitude to \emph{Cologero} for this generous and fair-minded review.

Publishing deadlines unfortunately prevented me from adequately treating the work of Eric Voegelin, which, upon reflection, I ought to have appealed to more frequently in the thesis under discussion here; his work, particularly \emph{The Political Religions} and \emph{The New Science of Politics} is highly recommended to readers.

Furthermore, I suggest that Voegelin be read in conjunction with Erik von Kuehnelt-Leddihn, particularly his \emph{Liberty or Equality}, \emph{Leftism Revisited} and \emph{Democracja -- Opium dla Ludu} (the later, as far as I am aware, is not yet translated into English). Thomas Molnar, too, is invaluable, specifically his \emph{The Counter-Revolution}.

These -- all of which feature in my bibliography -- will inoculate or at last provide an antidotum to the \emph{sinisterism} that often accompanies small-minded reading of Evola as well as the darker and unfortunately self-destructive tendencies of the ``alt-right''.

I am also grateful for having my attention drawn towards the work of Valentin Tomberg and Timothy Sprigge, whom I have only just discovered via \emph{Gornahoor}; thank you.

\hfill

\texttt{Mark Citadel on 2015-09-10 at 12:08 said:}

Linked and commented on this terrific article in my weekly essay:

\url{http://citadelfoundations.blogspot.com/2015/09/the-guiding-voice-of-old-right.html}


\hfill

\texttt{Cologero on 2015-09-10 at 13:18 said:}

Well, the credit goes to Edwin Diga for writing the original essay.

\hfill

\texttt{John Locke on 2015-09-18 at 12:18 said:}

``Nevertheless, he recognizes that there are differences between peoples, but that they are the result of group spiritual and psychical/cultural factors.''

50\% of the variance in all sorts of behavioral traits in humans are inherited (for some traits it's much higher) and most of the other 50\% cannot be explained by environmental factors, the cause is unknown. Man is not a perfectly malleable blank slate like some of the enlightenment philosophers and the soviets (`new Soviet man') believed. Nor is the soul entirely independent of the body. Man is a psychophysical whole, body and soul are interdependent and come into existence simultaneously. The soul is not absolutely transcendent with regard to the body, only the nous is.

This should be quite obvious, since brain development from childhood to adolescence and to maturity greatly alters the psyche. As does biochemistry.


\hfill

\end{sffamily}
\end{footnotesize}

